% !TeX program = lualatex
\documentclass{simplecv}
% \usepackage[utf8]{inputenc}
\usepackage{csquotes}
\usepackage[spanish, english]{babel}
\usepackage[margin=0.75in]{geometry}
\usepackage{dirtytalk}
\usepackage{multicol}
\usepackage[backend=biber,style=authoryear]{biblatex}
\usepackage{lipsum}

% Bibliography setup
\addbibresource{mybiblio.bib}

% Header information
\name{Roberto Olvera-Hernandez}
\title{Bioinformatics Software Engineer}
\email{roberto.olvhez@gmail.com}
\location{Zapopan}{Jalisco}{MX}
\phone{+52}{33-1540-9988}
\linkedin{\href{https://www.linkedin.com/in/rob-olvhez/}{rob-olvhez}}
\github{\href{https://github.com/robolvhez}{robolvhez}}

\newcommand{\mydot}{\raisebox{0.5 ex}{.}\ }

\begin{document}
\setmainfont{XCharter} % Sans-serif. Body.
\setsansfont{Cabin} % Sans. Titles.
\makecvheader{}
\thispagestyle{firstpage}

% ------------
% EXPERIENCE
% ------------
\cvsection{\job Professional \& Research Experience}

    \cvevent{Computational Genomics Researcher}
    {Center for Genomic Sciences (CCG) --- UNAM}
    {Cuernavaca, MX}
    {Feb. 2025}
    {Present}
    {
        \vspace{-5pt}
        Leading the design and implementation of scalable bioinformatics pipelines for processing genomic data from over 100,000 individuals on cloud infrastructure (DNAnexus; AWS/Azure).
        \begin{itemize}[leftmargin=*]
        \setlength{\itemsep}{0cm}
          \item \textbf{Software Development:} Engineered an ETL pipeline for large-scale genome-wide association studies (GWAS), optimizing data transformation and variant-calling workflows for high-throughput analyses.
          \item \textbf{Clinical data structring:} Designed and made a prototype of a OMOP data infrastructure for the prospective study, compending over 200 phenot.
          \item \textbf{AI-Assisted Engineering:} Leveraging AI coding tools for rapid prototyping, pipeline testing, and automation of complex data workflows.
          \item \textbf{Cross-functional Collaboration:} Partnering with international teams to ensure production-readiness and robust data governance for large-scale genetic pipelines.
        \end{itemize}
    }

    \cvevent{Genetic Epidemiology --- Research Placement}
    {Nuffield Department of Public Health (NDPH) --- University of Oxford}
    {Oxford, U.K.}
    {Aug. 2025}
    {Jan. 2026}
    {
        \vspace{-5pt}
        Developed and maintained a robust computational pipeline for within-family GWAS, directly addressing confounding from population structure in highly admixed cohorts.
        \begin{itemize}[leftmargin=*]
        \setlength{\itemsep}{0cm}
            \item \textbf{Data Pipelines:} Architected and deployed cloud-native pipelines for epidemiological data analysis, transitioning complex statistical models into scalable workflows.
            \item \textbf{Statistical Genetics:} Implemented advanced statistical genetics tools to extract unbiased heritability and genetic correlation estimates.
        \end{itemize}
    }

    \cvevent{Income Data Analyst}
    {TOKA Internacional}
    {Zapopan, MX}
    {2024}
    {2025}
    {
        \vspace{-5pt}
        Spearheaded the modernization of a fragile billing process for millions of monthly transactions, navigating a rapid migration from local servers to Azure.
        \begin{itemize}[leftmargin=*]
        \setlength{\itemsep}{0cm}
            \item \textbf{ETL \& Orchestration:} Architected and deployed an automated ETL pipeline using Python and Dagster, replacing fragile manual SQL processes and ensuring reliable fee processing.
            \item \textbf{Business Intelligence:} Automated the data delivery lifecycle, integrating the pipeline to publish results directly to Power BI and auto-email PDF/Excel reports to collaborators.
            \item \textbf{Cloud Migration:} Successfully adapted the data architecture and orchestration logic to Azure during a rapid cloud transition.
        \end{itemize}
    }

    \cvevent{Trainee}
    {Boehringer Ingelheim --- Vetmedica}
    {Zapopan, MX}
    {2021}
    {2022}
    {
        \vspace{-5pt}
        Bridged wet-lab operations and bioinformatics within an animal health laboratory, supporting viral genomic surveillance for veterinary vaccine development.
        \begin{itemize}[leftmargin=*]
        \setlength{\itemsep}{0cm}
            \item \textbf{Data Processing:} Engineered an automated data processing pipeline with Microsoft's Automation tools: extracting veterinary diagnostics information with AI from handwritten PDFs, directly into relational databases. Allowing collaborators to keep track of diagnosis through dashboards. 
            \item \textbf{Wet-Lab:} Executed Illumina MiSeq sequencing and qPCR, building a strong understanding of sample processing and DNA/RNA analysis.
            \item \textbf{Early bioinformatics:} Conducted preliminary viral evolution tracking across agricultural sites, analyzing genomic data to support animal health and vaccine research efforts.
        \end{itemize}
    }

% ------------
% EDUCATION
% ------------
\cvsection{\education Education}

    \cvevent{Ph.D. in Biomedical Sciences (Genomics Focus)}
    {Center for Genomic Sciences (CCG) --- UNAM}
    {Cuernavaca, MX}
    {2025}
    {Present}
    {}

    \cvevent{B.Sc. in Biotechnology Engineering}
    {ITESO}
    {Guadalajara, MX}
    {2018}
    {2023}
    {}

\vspace{0.3cm}
\begin{multicols}{2}
  
  % ------------
  % SKILLS
  % ------------
  \cvsection{\courses Technical Skills}
  \textbf{Programming:} Python, R, Bash, SQL, Git \\
  \textbf{Cloud/Infra:} DNAnexus, AWS (basics), Linux \\
  \textbf{Genomics:} Population \& Statistical Genetics, GWAS \\
  \textbf{AI Tools:} Agentic coding tools, Prompt engineering \\

  % ------------
  % LANGUAGES
  % ------------
  \cvsection{\lang Languages}
  \textbf{Spanish:} Native \\
  \textbf{English:} Fluent (TOEFL iBT - C1)

\vfill\null\columnbreak

  % ------------
  % CERTIFICATIONS
  % ------------
  \cvsection{\courses Certifications}
  \cvevent{Bioinformatic Tools for High-Throughput Sequencing}
      {USMBB - IBt UNAM}
      {}
      {}{Jul. 2023}
      {}
  \cvevent{Data Science: R Basics}
      {EBRC}
      {}
      {}{Aug. 2020}
      {}

\end{multicols}

\end{document}
